\setcounterpageref{ex}{0}
\TaoPhieuKiemTra{09}

\setcounter{ex}{0}
\setcounter{bt}{0}

\textbf{\large I. PHẦN TRẮC NGHIỆM (3,0 điểm)}

{\itshape Khoanh tròn vào chữ cái đặt trước câu trả lời đúng. Mỗi câu đúng được 0,5 điểm.}

\begin{ex} % Câu 1
Trong số $83$, chữ số $8$ có giá trị là
\choice
{$8$}
{\True $80$}
{$83$}
{$800$}
\loigiai{
Chữ số $8$ ở hàng chục nên có giá trị là $80$.\\
\textbf{Đáp án đúng: B}
}
\end{ex}

\begin{ex} % Câu 2
Trong phép tính $56 - 16 = 40$, số trừ là
\choice
{$56$}
{$40$}
{$72$}
{\True $16$}
\loigiai{
Số trừ là số $16$.\\
\textbf{Đáp án đúng: D}
}
\end{ex}

\begin{ex} % Câu 3
Thùng thứ nhất có $14\text{ }l$ nước, thùng thứ hai có $9\text{ }l$ nước. Thùng thứ nhất nhiều hơn thùng thứ hai là
\choice
{\True $5\text{ }l$}
{$4\text{ }l$}
{$6\text{ }l$}
{$23\text{ }l$}
\loigiai{
$14 - 9 = 5\text{ }l$.\\
\textbf{Đáp án đúng: A}
}
\end{ex}

\begin{ex} % Câu 4
Đoạn thẳng $CD$ dài $1\text{ dm}$. Đổi sang đơn vị xăng-ti-mét ta được
\choice
{$1\text{ cm}$}
{$100\text{ cm}$}
{\True $10\text{ cm}$}
{$20\text{ cm}$}
\loigiai{
$1\text{ dm} = 10\text{ cm}$.\\
\textbf{Đáp án đúng: C}
}
\end{ex}

\begin{ex} % Câu 5
So sánh: $11 - 4 \dots 12 - 6$.
\choice
{$<$}
{$=$}
{Không so sánh được}
{\True $>$}
\loigiai{
$11 - 4 = 7$; $12 - 6 = 6$. Vì $7 > 6$ nên điền dấu $>$.\\
\textbf{Đáp án đúng: D}
}
\end{ex}

\begin{ex} % Câu 6
Tính: $8 + 8 - 7 = \dots$
\choice
{$8$}
{\True $9$}
{$10$}
{$7$}
\loigiai{
$8 + 8 = 16$; $16 - 7 = 9$.\\
\textbf{Đáp án đúng: B}
}
\end{ex}

\vspace{4mm}
\textbf{\large II. PHẦN TỰ LUẬN (7,0 điểm)}

{\itshape Thí sinh trình bày chi tiết bài làm vào giấy kiểm tra.}

\begin{bt}[1,0 điểm]
Tính nhẩm:
\begin{multicols}{2}
\begin{enumerate}[a)]
  \item $9 + 8$
  \item $7 + 9$
  \item $16 - 6$
  \item $18 - 7$
\end{enumerate}
\end{multicols}
\loigiai{
\begin{itemize}
  \item $9 + 8 = 17$
  \item $7 + 9 = 16$
  \item $16 - 6 = 10$
  \item $18 - 7 = 11$
\end{itemize}
}
\end{bt}

\begin{bt}[1,5 điểm]
Đặt tính rồi tính:
\begin{multicols}{3}
\begin{enumerate}[a)]
  \item $41 + 34$
  \item $93 - 28$
  \item $59 + 38$
\end{enumerate}
\end{multicols}
\loigiai{
\begin{itemize}
  \item $41 + 34 = 75$
  \item $93 - 28 = 65$
  \item $59 + 38 = 97$
\end{itemize}
}
\end{bt}

\begin{bt}[1,0 điểm]
Tính giá trị biểu thức kèm đơn vị đo:
\begin{enumerate}[a)]
  \item $25\text{ kg} + 8\text{ kg} - 9\text{ kg}$
  \item $34\text{ }l - 15\text{ }l + 7\text{ }l$
\end{enumerate}
\loigiai{
\begin{enumerate}[a)]
  \item $25\text{ kg} + 8\text{ kg} - 9\text{ kg} = 33\text{ kg} - 9\text{ kg} = 24\text{ kg}$.
  \item $34\text{ }l - 15\text{ }l + 7\text{ }l = 19\text{ }l + 7\text{ }l = 26\text{ }l$.
\end{enumerate}
}
\end{bt}

\begin{bt}[1,5 điểm]
Bạn Tùng có $24$ quyển vở. Bạn Hoa có nhiều hơn bạn Tùng $7$ quyển vở. Hỏi bạn Hoa có bao nhiêu quyển vở?
\loigiai{
Số quyển vở bạn Hoa có là:
\[ 24 + 7 = 31 \text{ (quyển vở)} \]
\textbf{Đáp số:} $31$ quyển vở.
}
\end{bt}

\begin{bt}[1,0 điểm]
Một thùng chứa đầy $44\text{ }l$ nước mắm. Buổi sáng cửa hàng đã bán ra $13\text{ }l$. Hỏi trong thùng còn lại bao nhiêu lít nước mắm?
\loigiai{
Trong thùng còn lại số lít nước mắm là:
\[ 44 - 13 = 31\text{ }l \]
\textbf{Đáp số:} $31\text{ }l$ nước mắm.
}
\end{bt}

\begin{bt}[1,0 điểm]
Tìm số có hai chữ số, biết chữ số hàng chục là số lớn nhất có một chữ số, chữ số hàng đơn vị kém chữ số hàng chục $4$ đơn vị.
\loigiai{
Chữ số hàng chục là $9$.\\
Chữ số hàng đơn vị là: $9 - 4 = 5$.\\
Vậy số cần tìm là $95$.\\
\textbf{Đáp số:} $95$.
}
\end{bt}
